\section{Introduction}
\label{sec:intro}

Video generation models now produce photorealistic, temporally coherent footage, and on that basis they are increasingly described as world models that can reason about the visual world \citep{openai2024sora, deepmind2025veo3, wan2025wan, kuaishou2025kling, seedance2025}. VBVR-Pro puts that description to a test \citep{vbvrpro2026}. It asks whether a generative model can reason \emph{natively} in the visual medium: given a first frame and a prompt, the model must produce the video in which the prompt is carried out, and a rule-based evaluator checks the content of the frames. The suite is built from 300 procedural generators in five cognitive categories (abstraction, perception, spatiality, transformation, knowledge); its benchmark, VBVR-Pro-Bench, holds 100 tasks with five instances each, 50 in-domain families with training data and 50 out-of-domain families held out from all training. Video models score low on it: general-purpose systems land between 0.1 and 0.5, and the best, a Wan2.2 model RL-trained on the in-domain families with the evaluator as reward, reaches 0.670 overall and 0.532 out of domain.

We ask a different question about the same test. Every VBVR-Pro task is generated by a program, and its answer is a deterministic function of the first frame and the prompt. A model that denoises pixels must recover that function implicitly; a model that writes code can state it and let a renderer do the rest. So we swap the medium of the answer and nothing else. A coding agent is given the same first frame and the same prompt, together with the required container spec (frame size, frame rate, frame count, duration), and must write a program that renders the answer video. It sees no ground truth, metadata or generator, may not call an image or video generation model, runs in a CPU sandbox with common Python imaging libraries and a 30-minute limit, and gets one attempt per instance. The output goes to the unmodified official evaluator; an instance without a video scores zero, as it does for a video model that fails to generate (\cref{fig:teaser}). This is the program-synthesis route to visual reasoning \citep{gao2023pal, chen2023pot, suris2023vipergpt, gupta2023visprog}, applied to a benchmark built for video models and scored on their terms.

We ran this protocol on all 500 instances of VBVR-Pro-Bench with three closed-model agents (Codex with gpt-6-astra, Claude Code with Fable~5.1, Gemini CLI with Gemini~3.1~Pro) and, through the OpenCode harness, with 24 open-weight models. Video-model scores are the published VBVR-Pro leaderboard numbers, not rerun.

The gap is large (\cref{fig:leaderboard}, \cref{tab:main}). The three closed-model agents score 0.923, 0.894 and 0.893, all above 0.89, against 0.670 for the best video model. The RL-trained video model earns its score in domain (0.808) and loses it out of domain (0.532), and every trained video model falls on the held-out families. Every coding agent above the noise floor, closed or open, rises on them; Codex goes from 0.879 to 0.967 (\cref{fig:ood}). Among open-weight models the best, GLM-5, scores 0.572, above every video model except the RL-trained one; the bottom of the ranking is a failure of tool use, since models that never call a tool produce no video in 500 attempts and models that stop after one or two calls score near zero. The strongest agent is also cheap: Codex makes 3.2 tool calls per instance with a median wall-clock of 22 seconds (\cref{tab:efficiency}).

\paragraph{Contributions.}
\begin{enumerate}[leftmargin=*,itemsep=2pt,topsep=2pt]
\item \textbf{A medium swap, not a new benchmark.} A protocol under which a coding agent answers VBVR-Pro tasks by writing a program that renders the video, under the same inputs, output container and official evaluator as the video models.
\item \textbf{A leaderboard.} Three closed-model coding agents and 24 open-weight models scored on all 500 instances of VBVR-Pro-Bench, next to the published video-model leaderboard: 0.923 against 0.670.
\item \textbf{The out-of-domain reversal.} Trained video models lose their advantage on held-out task families and coding agents gain on them, reported per split and per cognitive category.
\item \textbf{An anatomy of agent behaviour.} Tool calls, wall-clock, tokens and failure modes per agent from the agents' own event logs, and the tasks on which even the strongest agents fail.
\end{enumerate}

We release the runner, sandbox and analysis code together with the per-instance evaluation results and run statistics behind every table in this paper at \url{https://github.com/hokindeng/solve-vision-with-code}; the project page is \url{https://video-reason.com}. Agent videos and run logs are available on request.
